% !TeX program = lualatex
% !TeX encoding = utf8
% !TeX root = ../Thermod.tex

%=========================================================
\Opensolutionfile{answer}[\currfilebase/\currfilebase-Answers]
\Writetofile{answer}{\protect\section*{\nameref*{\currfilebase}}}
\chapter{Метод термодинамічних потенціалів }\label{\currfilebase}
\makeatletter
\def\input@path{{\currfilebase/}}
\makeatother
%=========================================================

% --------------------------------------------------------
\section{1}
% --------------------------------------------------------
В термодинаміці розрізняють два методи: \emph{метод циклів} і \emph{метод термодинамічних потенціалів} або \emph{метод характеристичних функцій}. Перший метод вже використовувався раніше при виведенні деяких співвідношень. В методі послідовно обраховується:
\begin{enumerate}
	\item [а)] тепло, отримане системою в круговому процесі; 
	\item [б)] робота, виконана системою в циклі (вона або знаходиться з використанням геометричних міркувань або в подальшому використовується сам факт її існування);
	\item [в)] обраховується К.К.Д. циклу і зрівнюється з теоретичним.
\end{enumerate}

Цей метод принципово може бути застосований до будь-якої задачі, але має один недолік --- для установлення тої чи іншої закономірності кожен раз потрібно \emph{ad hoc} підбирати цикл: успішність вирішення залежить від вибору необхідного циклу, сам же вибір нічим не обумовлений. 

Вихідним в методі термодинамічних потенціалів є основне рівняння термодинаміки:
\begin{equation}\label{5.1}
	TdS = dU+\sum_{i}F_idx_i
\end{equation}
яке дозволяє ввести різні функції стану, адекватні відповідним умовам --- термодинамічні потенціали. Для визначеності розглянемо систему з параметрами $T$, $P$, $V$:
\begin{equation}\label{5.2}
	TdS = dU+PdV
\end{equation}
Рівняння (\ref{5.2}) пов'язує п'ять функцій стану, дві з яких повністю визначають стан простої системи. Для знаходження інших можна скористатися термічним і калоричним рівняннями стану: $P=P(V,T)$ і $U=U(V,T)$. Але при належному виборі незалежних змінних за певних умов можна використовувати лише одне додаткове рівняння замість рівнянь стану. Запишемо (\ref{5.2}) в вигляді
\begin{equation}\label{5.3}
	dU=TdS-PdV
\end{equation}
і обираючи як незалежні змінні $S$ і $V$, переконаємося, що достатньо знати $U$ як функцію ентропії і об'єму $U=U(S,V)$. В силу того, що приріст внутрішньої енергії $dU$ є повний диференціал, диференціюванням $U$ знайдемо $Т$ і $P$:
\begin{equation}\label{5.4}
	T=\lpdr{U}{S}{V}; ~~~~ P=-\lpdr{U}{V}{S}
\end{equation}
Повторне диференціювання дозволяє знайти термодинамічні коефіцієнти:
\begin{equation}\label{5.5}
	\dfrac{T}{C_V}=\left(\dfrac{\partial^2U}{\partial S^2}\right)_{V}; ~~~~ K_{S}=\gamma_{S}^{-1}=\left(\dfrac{\partial^2U}{\partial V^2}\right)_{S},
\end{equation}
а змішані похідні, рівність яких випливає з того, що $dU$ є повний диференціал, дають нам одне із співвідношень Максвела (\ref{4.20}):
\begin{equation}\label{5.6}
	\dfrac{\partial^2U}{\partial V\partial S} = \lpdr{T}{V}{S} = - \lpdr{P}{S}{V}=\dfrac{\partial^2U}{\partial V\partial S}
\end{equation}
Функцію $U$ називають \emph{термодинамічним потенціалом}, а змінні $V$ i $S$ \emph{натуральними} або \emph{природними змінними}. Пов'язано це з тим, що подібно до сил в механіці узагальнені сили $P$ і $T$ визначаються диференціюванням по узагальнені координатам. Внутрішню енергію також називають \emph{адіабатним потенціалом}, оскільки робота зовнішніх сил дорівнює із зворотним знаком зміні внутрішньої енергії в адіабатичних процесах. Внутрішня енергія, визначена в змінних $S$ і $V$, є \emph{характеристичною функцією}, оскільки всі інші функції стану визначаються диференціюванням її по $S$ і $V$.
% --------------------------------------------------------
\section{2}
% --------------------------------------------------------